\begin{figure}[tb]
\centering
\begin{tikzpicture}[x=0.97cm,y=0.92cm,font=\small]
 \filldraw[fill=blue!12,draw=blue!65!black] (0,0) rectangle (4.8,1.05);
 \node at (2.4,.525) {$W_e$: $r_h$ indices};
 \draw[gray] (4.8,0) rectangle (5.4,1.05);
 \node at (5.1,.525) {$g$};
 \filldraw[fill=teal!12,draw=teal!60!black] (5.4,0) rectangle (9.3,1.05);
 \node[align=center] at (7.35,.525) {child-subtree windows\\and gaps: $\sigma_{h-1}$};
 \draw[gray] (9.3,0) rectangle (9.9,1.05);
 \node at (9.6,.525) {$g$};
 \filldraw[fill=black!5,draw=gray] (9.9,0) rectangle (12.3,1.05);
 \node[align=center] at (11.1,.525) {next sibling\\block};
 \node[above] at (0,1.1) {$a_e$};
 \node[above] at (4.8,1.1) {$b_e$};
 \node[above] at (9.3,1.1) {$b_e^*$};
 \draw[decorate,decoration={brace,mirror,amplitude=5pt},thick]
   (0,-.16)--(9.3,-.16);
 \node[below,align=center] at (4.65,-.42)
   {$b_e^*-a_e+1=r_h+g+\sigma_{h-1}\le2r_h$};
 \draw[-{Stealth[length=2mm]},gray] (-.2,-1.25)--(12.5,-1.25);
 \node[below] at (6.15,-1.25) {increasing indices in preorder};
\end{tikzpicture}
\caption{The block attached to an edge at height $h\ge2$, shown
schematically. The window length is at least the length of the following
gap and child-subtree span together. Thus the entire block has
at most twice the window length, regardless of where it occurs in the
tree. The next sibling block is outside this bound. At height $1$,
the child is a leaf and the block consists only of $W_e$.}
\label{fig:window-block}
\end{figure}
